%% Describe what is the problem current algos (stability isses, popularity influence): Gradient Ascent, Gradient Difference, NPO, WGA;
%% Explain how we solve this popularity gap 
%% Which benchmark we use to evaluate (we use real-world facts QA benchmarks for solving real-words tasks; we scoring facts based on they popularity. if you have ascess to pretrain set, you can use information from it to scoring populaeiry; otherwise use 'approximate' metrics from large data corpus like wikidata to score facts. here popularity is the global metric for for unlerning regularization; the local (model) metric is log probes and its to factors influence to unlearning.) also for trade-off stability we automatically balance retain, that directly effect to forgetting.

The necessity of removing specific knowledge from Large Language Models (LLMs) has led to the development of various weight-updating unlearning algorithms. However, established methods often struggle with a fundamental stability trade-off between forgetting target information and preserving general model capabilities. Traditional Gradient Ascent (GA) simply maximizes the loss on the forget set, which frequently leads to catastrophic collapse of the model's linguistic coherence. Gradient Difference (GD) attempts to mitigate this by adding a retain loss, yet it still treats every token in the forget set as having equal importance, regardless of whether a token is a critical factual entity or a high-frequency noise word like an article or preposition. More recent approaches such as Negative Preference Optimization (NPO) and Weighted Gradient Ascent (WGA) attempt to refine the unlearning process by using the model's own internal states. NPO treats unlearning as a preference alignment problem where the model is pushed away from the "preferred" response. WGA uses the model's log-probabilities to weight the gradient updates, theoretically focusing the unlearning effort on tokens where the model is most confident. While these methods are more sophisticated than simple ascent, they rely exclusively on the model's current confidence as a local metric. Our core assumption is that log-probabilities alone do not always accurately reflect how deeply a fact is embedded in the model. We argue that incorporating an external proxy for "fact popularity" better conveys the necessary unlearning strength than internal model confidence alone.

To address this, we introduce AdaPop (Adaptive Popularity), an algorithm designed to make the unlearning process more flexible by modulating updates based on the global prevalence of a fact. The relationship between a fact's presence in a massive pretraining corpus and its resistance to unlearning is often non-linear. Popularity can be quantified directly if the pretraining data is accessible by measuring citelinks or co-occurrence frequencies. In cases where the original pretraining corpus is private or unknown, we utilize large-scale public data proxies. This includes structural knowledge bases like Wikidata or massive web-scale corpora like Common Crawl to estimate how often a subject-object relationship appears in the general distribution of human knowledge.

AdaPop uses a popularity-dependent exponent to adjust the influence of each fact during the optimization process. This ensures that the gradient updates are scaled appropriately: deeply entrenched, popular facts are handled with a different weighting logic than obscure, long-tail information that might only appear a handful of times in the training set. Furthermore, AdaPop implements a dual-ascent controller to automatically manage the balance between forgetting and retaining. This removes the need for manual hyperparameter tuning of the retain coefficient $\alpha$, allowing the algorithm to adaptively maintain model utility based on real-time drift statistics. We demonstrate the effectiveness of this popularity-aware approach on benchmarks like DUET and RWKU, where diverse fact distributions highlight the limitations of purely confidence-based methods.